% !TEX root = ../thesis.tex

\chapter{Related Work} \label{chp:relwork}

\section{RL-Based Legged Locomotion}

Reinforcement learning has become the dominant paradigm for quadruped locomotion
over the past five years.
Lee~et~al.~\cite{lee2020learning} demonstrated that a neural locomotion policy
trained entirely in simulation can transfer to real quadrupeds and traverse
challenging terrain including stairs, sloped ground, and loose gravel without
additional fine-tuning, establishing simulation-to-reality transfer as a
viable path.
Kumar~et~al.~\cite{kumar2021rma} proposed Rapid Motor Adaptation (RMA), a
two-module architecture in which a base policy trained with privileged
simulation information is paired with a fast adaptation module that estimates
environment parameters from recent proprioceptive history; the combined system
recovers from perturbations in real time.
Rudin~et~al.~\cite{rudin2022learning} showed that massively parallel simulation
in Isaac Gym~\cite{makoviychuk2021isaac} reduces wall-clock training time to
minutes, making iterative reward shaping practical.
Margolis and Agrawal~\cite{margolis2022walk} introduced a command-conditioned
locomotion policy (\textit{Walk These Ways}) capable of expressing diverse gaits
including trot, bound, and pronk through a single unified policy conditioned on
gait parameters such as frequency, phase offset, and foot-swing height.
This distributional command representation is directly adopted by the RoboDuet
framework~\cite{pan2024roboduet} and, by inheritance, by our B2+Z1 implementation.
Cheng~et~al.~\cite{cheng2024parkour} pushed the boundary further with
\emph{Extreme Parkour}, training a policy to jump, climb, and slide over
obstacles using only proprioception and a single depth camera, demonstrating
the expressive capacity of PPO-based policies on complex agility tasks.

\section{Legged Loco-Manipulation}

Equipping a legged robot with a manipulator arm introduces the
\textit{loco-manipulation} challenge: the arm and legs must be coordinated to
prevent the arm's inertial loads from destabilising locomotion, and the body
posture must be adjusted to extend the arm's reachable workspace.

Early work addressed this via whole-body model predictive control (MPC).
Sleiman~et~al.~\cite{sleiman2021unified} presented a unified MPC framework for
ANYmal~\cite{hutter2016anymal} with a fixed-base arm, enabling simultaneous
walking and end-effector tracking.
While model-based methods offer formal guarantees and high-frequency control,
they demand accurate system models and struggle with unknown contacts or
rapidly changing environments.

Learning-based methods have begun to close this gap.
Ma~et~al.~\cite{ma2022combining} combined a learned locomotion policy with a
model-based arm controller, demonstrating that a decoupled architecture can
achieve basic pick-and-place on a legged platform.
Fu~et~al.~\cite{fu2023deep} trained a \emph{Deep Whole-Body Control} policy
for a quadruped arm system using a teacher--student distillation scheme,
achieving dynamic arm motions during locomotion but requiring separate training
of a privileged teacher and a deployable student.
Jeon~et~al.~\cite{jeon2023learning} applied whole-body RL to a quadrupedal
robot to track arbitrary end-effector trajectories, using a single monolithic
policy that outputs both leg and arm joint targets simultaneously.

\section{RoboDuet: Cooperative Whole-Body Loco-Manipulation}

The most directly relevant prior work is RoboDuet by Pan~et~al.~\cite{pan2024roboduet},
which is the base framework reproduced and extended in this project.
RoboDuet departs from monolithic policy designs by introducing \emph{two
cooperating policies}: a \textbf{locomotion (dog) policy} responsible for body
movement and a separate \textbf{arm policy} responsible for end-effector
control.
The key architectural innovation is the \emph{communication channel}: the arm
policy outputs not only joint position targets for the six arm joints, but also
a desired body-orientation command (pitch and roll) that is \emph{injected} into
the observation vector of the locomotion policy, replacing the user-specified
body-orientation command.
This allows the arm policy to dynamically lean the body to expand the reachable
workspace without any explicit coordination signal from the user.

Both policies are trained with PPO using separate actor-critic networks and
separate experience buffers, but they share the same simulation environment.
Training proceeds in two stages: first the locomotion policy is pre-trained
until stable gait is achieved, then the arm policy is introduced and both
policies are fine-tuned jointly.

The original RoboDuet experiments used the Unitree \textbf{Go1} quadruped
(body mass $\approx\!12$\,kg) with the Unitree Z1 arm
(6-DoF, payload $\approx\!2$\,kg).
Our work adapts this framework to the Unitree \textbf{B2} quadruped
(body mass $\approx\!52$\,kg), which requires non-trivial changes to the
PD servo gains and training curriculum to compensate for the significantly
different mass distribution and joint torque characteristics,
as detailed in Chapter~\ref{chp:method}.

\section{Summary and Positioning}

Table~\ref{tab:relwork} summarises the key related works along the dimensions
most relevant to this project.
Our work occupies the intersection of RL-based locomotion, arm manipulation,
and cooperative dual-policy architecture, on a platform (B2) that has not
previously been demonstrated in a loco-manipulation setting.

\begin{table}[h]
\centering
\caption{Comparison of related work. ``Loco-manip.'' denotes simultaneous
locomotion and arm manipulation. ``Dual policy'' denotes a separate
locomotion and arm policy with inter-policy communication.}
\label{tab:relwork}
\small
\begin{tabular}{lcccc}
\hline
\textbf{Work} & \textbf{RL-based} & \textbf{Loco-manip.} & \textbf{Dual policy} & \textbf{Platform} \\
\hline
Lee et al.~\cite{lee2020learning}       & \checkmark & --         & --         & ANYmal \\
Margolis \& Agrawal~\cite{margolis2022walk} & \checkmark & --         & --         & A1     \\
Sleiman et al.~\cite{sleiman2021unified} & --         & \checkmark & --         & ANYmal \\
Fu et al.~\cite{fu2023deep}              & \checkmark & \checkmark & --         & A1+arm \\
Jeon et al.~\cite{jeon2023learning}      & \checkmark & \checkmark & --         & custom \\
Pan et al.~\cite{pan2024roboduet}        & \checkmark & \checkmark & \checkmark & Go1+Z1 \\
\textbf{Ours}                            & \checkmark & \checkmark & \checkmark & \textbf{B2+Z1} \\
\hline
\end{tabular}
\end{table}
